\section{Identifying the cyclic cubic quotient}\label{sec:quotient}

Let $L$ be the splitting field of $f$ over $\Q(t)$ and put $M=L^N$.
By \cref{prop:certificate}, inertia at the two finite branch points has
order $3$, and inertia at infinity has order $13$.
\Cref{lem:rationalquotient} gives $M=\Q(s)$. To obtain an explicit
polynomial with group $N$, we determine $t$ as a rational function of $s$.

\subsection{The fiber at infinity}\label{sec:infinity}

Fix a place of $L$ above $\infty$, with inertia $I$ and decomposition
group $D$. Tame ramification of index $13$ puts $\mu_{13}$ in its residue
field $k$. Since the base place is rational,
\[
 D/I=\Gal(k/\Q)\twoheadrightarrow
 \Gal(\Q(\zeta_{13})/\Q)\simeq C_{12}.
\]
On the other hand, $D\subseteq N_G(I)$ and the latter has order $156$.
Thus $D=N_G(I)$ and $k=\Q(\zeta_{13})$.
Moreover, $D\cap N$ has order $52$. The residue field in the cubic
quotient is therefore the unique cubic subfield of $\Q(\zeta_{13})$,
represented by
\[
 z^3+z^2-4z+1.
\]
The orbits of $D$ on the $65$ letters have lengths $13$ and $52$.
The inertia group has five orbits of length $13$; $D/I$ fixes one of
these orbits and permutes the other four transitively. The corresponding
closed points therefore have degrees $13/13=1$ and $52/13=4$.
Denote these points by $P_0$ and $P_4$, respectively. On the degree-$65$
source curve the pole divisor is therefore
\begin{equation}\label{eq:polefiber}
 (t)_\infty=13(P_0+P_4).
\end{equation}
The residue field at $P_4$ is the unique quartic subfield of
$\Q(\zeta_{13})$. Geometrically, $P_4$ becomes four points, each of
ramification index $13$; together they account for the $52$ sheets,
while $P_0$ accounts for the other $13$. See also \cref{fig:tower}.

\subsection{Two twists and one distinguishing prime}

Consider the reference cyclic cubic cover
\begin{equation}\label{eq:shanks}
 u^3-3tu^2-(3t+3)u-1=0.
\end{equation}
Its discriminant is $81(t^2+t+1)^2$, its deck group is constant $C_3$,
and its fiber at infinity consists of three rational points. Over
$\overline{\Q}$, a cyclic cubic cover with these two branch points is
unique up to isomorphism over the base. Its $\Q$-forms are consequently
classified by
\[
 H^1(\Q,C_3)=\operatorname{Hom}_{\mathrm{cont}}
 (\Gal(\overline{\Q}/\Q),C_3),
\]
because the deck group is constant.  The split fiber over infinity is a
$C_3$-torsor.  Twisting by a character $\chi$ changes its Galois action
from the trivial action to translation by $\chi$; when $\chi$ is
surjective, its residue field is therefore the fixed field of
$\ker(\chi)$.  Requiring this field to be the conductor-$13$ cubic field
$K$ leaves precisely the two isomorphisms
$\Gal(K/\Q)\simeq C_3$, which are inverse to one another.  One can
construct their twists over
the cyclic cubic field $K=\Q[z]/(z^3+z^2-4z+1)$: if $\tau$ generates
$\Gal(K/\Q)$ and $\rho(u)=-1/(u+1)$ generates the deck group of
\eqref{eq:shanks}, descend $K(u)$ by the semilinear actions
$(\tau,\rho)$ and $(\tau,\rho^{-1})$. Rational coordinates on the two
fixed fields give the following maps, whose discriminants and fibers
provide exact checks of the construction:
\begin{align}
 T_1(s)&=\frac{-4s^3+6s^2+36s-6}{3(s^3+5s^2-9s-5)},\label{eq:T1}\\
 T_2(s)&=\frac{s^3+6s^2-27s+21}{3(s^3-7s^2+12s-5)}.\label{eq:T2}
\end{align}
Writing $T_i=n_i/d_i$ with the displayed numerators and denominators,
the exact identities
\[
 \disc_s(n_1-td_1)=936^2(t^2+t+1)^2,\qquad
 \disc_s(n_2-td_2)=117^2(t^2+t+1)^2
\]
verify their branch loci and cyclicity. Both denominator fields are
isomorphic to the conductor-$13$ field above. Their fibers at $t=0$
are the distinct cyclic cubic fields represented by $z^3-39z-26$ and
$z^3-39z-91$, respectively. Both have field discriminant $117^2$, hence
conductor $117=9\cdot13$: they combine the conductor-$9$ fiber of the
reference cover at $t=0$ with the two conductor-$13$ characters.
The accompanying program
\texttt{scripts/check\_cubic\_maps.py} checks these discriminants and
field isomorphisms exactly with PARI/GP.

At a prime of separable reduction, the fiber of $M$ over $t=0$ splits
if and only if the Frobenius element of $f(X,0)$ belongs to $N$.
Every element of $G\setminus N$ has order divisible by $3$, and $N$ has
no element of order $3$. Thus membership is read from the factor degrees:
the Frobenius belongs to $N$ precisely when none is divisible by $3$.

At $p=29$, the factorization type of $f(X,0)$ is $1\,2^2\,6^{10}$, so
Frobenius lies outside $N$. The cubic for $T_1$ splits modulo $29$,
whereas the cubic for $T_2$ does not. Since the two twists exhaust the
possibilities, this single splitting test identifies $M$ with $\Q(s)$
via $t=T_2(s)$.

The equality $L^N=\Q(s)$ proves that $L/\Q(s)$ is regular with group
$N$. Since $N$ is transitive on the $65$ roots, $f(X,T_2(s))$ is
irreducible. For a polynomial with coefficients in $\Q[s]$ we clear
denominators:
\begin{equation}\label{eq:g}
 g(X,s)=\bigl(3(s^3-7s^2+12s-5)\bigr)^7 f(X,T_2(s)).
\end{equation}
The polynomial $g$ has bidegree $(65,21)$ and the same splitting field.
Its leading coefficient in $X$ is the displayed seventh power; the
monic polynomial is $f(X,T_2(s))\in\Q(s)[X]$.
